\section{Módulo 4: Variabilidad en las Operaciones}

\subsection{Por qué importa la variabilidad}

\begin{definicion}{El problema central}
La causa de las colas es que, en el corto plazo y en períodos impredecibles, la \textbf{capacidad no logra satisfacer la demanda}. Hay que distinguir \textbf{variabilidad} (predecible, regular) de \textbf{incertidumbre}. El trade-off de las colas: el costo de mejor servicio sube, mientras el costo de espera del cliente baja; el óptimo cruza ambas curvas.
\end{definicion}

\begin{formulas}{Ley de Little y su extrapolación productiva}
\[
  L = \lambda \times W \qquad\Longrightarrow\qquad \text{WIP} = \text{TH} \times \text{CT}
\]
$L$ = inventario en el sistema, $\lambda$ = tasa de llegada, $W$ = tiempo en el sistema. En la fábrica: WIP = trabajo en proceso, TH = throughput, CT = tiempo de ciclo. Vale en \textbf{cualquier} disciplina de cola en \textbf{estado estacionario} (J.D.\ Little, MIT, 1961). Para reducir CT: bajar WIP o subir TH.
\end{formulas}

\subsection{Medir la variabilidad}

\begin{formulas}{Coeficiente de variación}
\[
  c = \frac{\sigma}{t} = \frac{\sqrt{\text{Var}(T)}}{E(T)}
\]
Clasificación: \textbf{Baja} ($c<0{,}75$), \textbf{Moderada} ($0{,}75 \le c \le 1{,}33$), \textbf{Alta} ($c>1{,}33$). Es la \textbf{mejor forma} de medir variabilidad (adimensional).
\end{formulas}

\begin{formulas}{Variabilidad por fallas}
\begin{align*}
  A &= \frac{m_f}{m_f + m_r} && \text{(disponibilidad)}\\
  r_e &= A\,r_0, \qquad t_e = \frac{t_0}{A} && \text{(tasa y tiempo efectivo)}\\
  c_e^2 &= c_0^2 + (1 + c_r^2)\,A(1-A)\,\frac{m_r}{t_0} && \text{(CV efectivo)}
\end{align*}
$m_f$ = tiempo medio entre fallas, $m_r$ = tiempo medio de reparación, $c_r$ = CV de la reparación. \textbf{Conclusión:} paradas \textbf{largas e infrecuentes} aumentan mucho más $c_e^2$ que paradas cortas y frecuentes.
\end{formulas}

\begin{formulas}{Variabilidad por setups}
\[
  t_e = t_0 + \frac{t_s}{N_s}, \qquad
  \sigma_e^2 = \sigma_0^2 + \frac{\sigma_s^2}{N_s} + \frac{N_s - 1}{N_s^2}\,t_s^2, \qquad
  c_e^2 = \frac{\sigma_e^2}{t_e^2}
\]
$t_s$ = tiempo de setup, $N_s$ = unidades entre setups. El setup se prorratea sobre el lote.
\end{formulas}

\begin{formulas}{Propagación de la variabilidad (látigo)}
\[
  (c_s)^2 \approx \rho^2 (c_e)^2 + (1 - \rho^2)(c_a)^2
\]
La variabilidad de \textbf{salida} ($c_s$) depende de la variabilidad de servicio ($c_e$), la de llegadas ($c_a$) y la saturación ($\rho$). ``La variabilidad se propaga (efecto látigo)''.
\end{formulas}

\subsection{Modelos de colas}

\begin{formulas}{Sistema M/M/1}
Llegadas Poisson, servicio exponencial, 1 servidor. Estable solo si $\rho<1$.
\begin{align*}
  \rho &= \frac{\lambda}{\mu} && \text{(utilización)}\\
  L &= \frac{\rho}{1-\rho}, \qquad W = \frac{1}{\mu(1-\rho)} && \text{(en el sistema)}\\
  L_q &= \frac{\rho^2}{1-\rho}, \qquad W_q = \frac{\rho}{\mu(1-\rho)} && \text{(en la cola)}
\end{align*}
\end{formulas}

\begin{formulas}{Sistema G/G/1 --- Ecuación de Kingman (VUT)}
Para distribuciones generales, el tiempo en cola es:
\[
  CT_q = \underbrace{\left(\frac{c_a^2 + c_e^2}{2}\right)}_{V\ (\text{Variabilidad})}
         \underbrace{\left(\frac{\rho}{1-\rho}\right)}_{U\ (\text{Utilización})}
         \underbrace{\vphantom{\frac{\rho}{1-\rho}}\ t_e\ }_{T\ (\text{Tiempo})}
\]
\textbf{Kingman: CT $= V\cdot U\cdot T$}. Es exacta para M/M/1 (cuando $c_a=c_e=1$). La espera es el producto de \textbf{Variabilidad, Utilización y Tiempo}.
\end{formulas}

\begin{formulas}{Sistemas restringidos: M/M/1/b y M/M/c}
\textbf{Capacidad finita (M/M/1/b)}, cola con tope $b$, hay bloqueo:
\[
  L = \frac{\rho}{1-\rho} - \frac{(b+1)\rho^{b+1}}{1-\rho^{b+1}}, \qquad
  \lambda' = \lambda\left(\frac{1-\rho^b}{1-\rho^{b+1}}\right)
\]
\textbf{Múltiples servidores (M/M/c)}, $\rho = \dfrac{\lambda}{c\mu}$:
\[
  L_q = \frac{\rho}{1-\rho}\times \text{Prob}(N>c), \qquad
  W_q = \frac{\rho}{\lambda(1-\rho)}\times \text{Prob}(N>c)
\]
\end{formulas}

\begin{trampa}{La utilización explota la cola cerca de $\rho=1$}
A medida que $\rho \to 1$, el tiempo en cola \textbf{explota exponencialmente}. Un pequeño cambio en la utilización (de 0,90 a 0,95) genera una diferencia drástica. Por eso, con alta utilización se requieren \textbf{grandes} reducciones de variabilidad (motivación del 6-Sigma). Nunca planifiques operar al 100\%.
\end{trampa}


\subsection[Variabilidad: Optimización G/G/1]{Ejercicio de Aplicación: Optimización de Variabilidad y Tasa de Llegada G/G/1 (Ayudantía 8)}

\begin{ejercicio}{Modelo de cola G/G/1 con elasticidad de demanda y capacidad (Ayudantía 8)}
\begin{tipprueba}{¿Cuándo usar este enfoque?}
\textbf{Identificación:} te dan una cola G/G/1 con tasa de llegada elástica respecto a un descuento $\Delta$ ($\lambda = \lambda_0 e^{-\Delta}$), costos de espera lineales y tiempos de servicio generales. Piden plantear el modelo matemático de optimización y sus condiciones de primer orden (CPO) para hallar el descuento óptimo.
\textbf{Qué aplicar:} ecuación de Kingman para el tiempo de espera, modelar la función de costo total (descuento acumulado + costo de espera del cliente), y derivar con respecto a la variable de decisión.
\end{tipprueba}

\textbf{Enunciado:}
Los clientes de una imprenta llegan a una tasa $\lambda$ [clientes/hr] (distribución general) y la capacidad de atención es $\mu$ [clientes/hr]. El tiempo efectivo de procesamiento es $t_e$, el coeficiente de variabilidad de llegadas es $c_a$ y el de servicio es $c_e$. Para evitar colapsos y setups costosos, se evalúa realizar descuentos ($\Delta$) en el precio de venta para alterar la tasa de llegada, de modo que:
$$\lambda = \lambda_0 e^{-\Delta}$$
Los clientes tienen una función de costo de espera total $CE(W) = 100 + W$, donde $W$ es el tiempo total en el sistema (espera en cola más servicio).
1. Plantee el modelo de programación matemática que permita optimizar el proceso productivo. Deje expresadas las condiciones de primer orden (CPO) que permitan obtener el óptimo, sin resolverlas.
2. Si se abre la posibilidad de adquirir nueva capacidad productiva y aumentar su tasa de atención a un costo de \$K por cada unidad de aumento, ¿cómo cambia el modelo?

\textbf{Solución:}

\textbf{1) Formulación del Modelo de Optimización:}
\begin{itemize}
  \item \textbf{Variable de decisión:} Descuento otorgado por cliente, $\Delta$.
  \item \textbf{Tasa de llegada elástica:} $\lambda(\Delta) = \lambda_0 e^{-\Delta}$.
  \item \textbf{Tiempo total en el sistema ($W$):}
\end{itemize}
    Por la ecuación de Kingman para colas $G/G/1$:
    $$W = CT_q + t_e = \left( \frac{c_a^2 + c_e^2}{2} \right) \left( \frac{\rho}{1 - \rho} \right) \frac{1}{\mu} + \frac{1}{\mu}$$
    Sustituyendo $\rho = \lambda(\Delta) / \mu$ y $t_e = 1/\mu$:
    $$W(\Delta) = \left( \frac{c_a^2 + c_e^2}{2} \right) \left( \frac{\lambda_0 e^{-\Delta}}{\mu(\mu - \lambda_0 e^{-\Delta})} \right) + \frac{1}{\mu}$$
\begin{itemize}
  \item \textbf{Función de costo total a minimizar ($f(\Delta)$):}
\end{itemize}
    $$\min_{\Delta} \quad f(\Delta) = \lambda(\Delta) \cdot \Delta + \lambda(\Delta) \cdot CE(W(\Delta))$$
    $$\min_{\Delta} \quad f(\Delta) = \lambda_0 e^{-\Delta} \left[ \Delta + 100 + \left( \frac{c_a^2 + c_e^2}{2} \right) \left( \frac{\lambda_0 e^{-\Delta}}{\mu(\mu - \lambda_0 e^{-\Delta})} \right) + \frac{1}{\mu} \right]$$
\begin{itemize}
  \item \textbf{Restricciones:}
  \item \textit{Estabilidad del sistema:} $\rho < 1 \implies \lambda_0 e^{-\Delta} < \mu \implies \Delta > \ln\left( \frac{\lambda_0}{\mu} \right)$
  \item \textit{No negatividad:} $\Delta \ge 0$
\end{itemize}

\begin{itemize}
  \item \textbf{Condiciones de Primer Orden (CPO):}
\end{itemize}
    Derivando la función objetivo $f(\Delta)$ con respecto a $\Delta$ e igualando a cero:
    $$\frac{df}{d\Delta} = 0 \implies -\lambda(\Delta) \left[ \Delta + 100 + W(\Delta) \right] + \lambda(\Delta) \left[ 1 + \frac{dW}{d\Delta} \right] = 0$$
    Simplificando por $\lambda(\Delta)$ (dado que $\lambda(\Delta) > 0$):
    $$W(\Delta) + 100 + \Delta - 1 - \frac{dW}{d\Delta} = 0$$
    Donde la derivada del tiempo de espera es:
    $$\frac{dW}{d\Delta} = \left( \frac{c_a^2 + c_e^2}{2} \right) \left( \frac{-\lambda_0 e^{-\Delta}}{\mu - \lambda_0 e^{-\Delta}} - \frac{\lambda_0^2 e^{-2\Delta}}{(\mu - \lambda_0 e^{-\Delta})^2} \right) \frac{1}{\mu}$$

\textbf{2) Modelo con Ampliación de Capacidad:}
\begin{itemize}
  \item \textbf{Nuevas variables de decisión:} Descuento $\Delta \ge 0$ y ampliación de capacidad $u \ge 0$.
  \item La nueva tasa de servicio es $\mu = \mu_0 + u$.
  \item La función objetivo ahora incluye el costo de expansión de capacidad $K \cdot u$:
\end{itemize}
    $$\min_{\Delta, u} \quad f(\Delta, u) = \lambda_0 e^{-\Delta} \left[ \Delta + 100 + W(\Delta, u) \right] + K \cdot u$$
    Donde:
    $$W(\Delta, u) = \left( \frac{c_a^2 + c_e^2}{2} \right) \left( \frac{\lambda_0 e^{-\Delta}}{(\mu_0 + u)(\mu_0 + u - \lambda_0 e^{-\Delta})} \right) + \frac{1}{\mu_0 + u}$$
\begin{itemize}
  \item \textbf{Restricciones de estabilidad:}
\end{itemize}
    $$\lambda_0 e^{-\Delta} < \mu_0 + u$$
    $$\Delta \ge 0, \quad u \ge 0$$
\end{ejercicio}

\subsection{Estrategias y psicología}

\begin{definicion}{Pooling, buffers y psicología de la espera}
\begin{itemize}[leftmargin=1.4em, itemsep=1pt]
  \item \textbf{Pooling (agrupación):} consolidar recursos (unifila, inventario centralizado) promedia y \textbf{disminuye} la variabilidad global.
  \item \textbf{Buffer finito:} amortigua la propagación de variabilidad, pero \textbf{reduce el throughput} máximo. Limitar el buffer reduce CT (Little).
  \item \textbf{Psicología:} el tiempo ocioso se percibe más largo; las esperas inciertas, inexplicadas e injustas parecen más largas; las esperas individuales más largas que las grupales. Recomendaciones: distraer al cliente, informar la espera, ocultar empleados inactivos, segmentar.
\end{itemize}
\end{definicion}

\subsection{V/F frecuentes de Variabilidad}

\begin{center}
\small
\begin{tabular}{@{}clp{10.5cm}@{}}
  \toprule
  $\#$ & R & Afirmación y corrección \\
  \midrule
  1 & \textbf{V} & ``La ecuación de Kingman muestra que la espera es el producto de Variabilidad, Utilización y Tiempo.'' \\
  2 & \textbf{F} & ``Un buffer de capacidad finita aumenta el throughput del sistema.'' --- Falso: lo \textbf{reduce}, aunque mejore el servicio a quienes ingresan. \\
  3 & \textbf{F} & ``Paradas cortas y frecuentes aumentan más la variabilidad que paradas largas e infrecuentes.'' --- Falso: es al revés. \\
  4 & \textbf{V} & ``La Ley de Little es válida para cualquier disciplina de cola en estado estacionario.'' \\
  5 & \textbf{F} & ``Aumentar la utilización siempre mejora el desempeño.'' --- Falso: cerca de $\rho=1$ la cola explota exponencialmente. \\
  6 & \textbf{V} & ``El pooling de recursos reduce la variabilidad global del sistema.'' \\
  \bottomrule
\end{tabular}
\end{center}

% ════════════════════════════════════════════════════════════════════════════


